% LAB ONLY: existing render_block with canonical-row mapping; not full legacy renderer or admission
\begin{claimboundary}
Mistral Agentic SearchのFinanceBench／OfficeQA Proの数値は、既定の構成を使ったMistral自身の報告である。今回確認された資料には独立した再現検証がなく、提供範囲や料金についても発表記事を超える詳細は確認されていない。\autocite{src-a9b696153570}
\end{claimboundary}
\begin{claimboundary}
Slack Codeの企業向け展開や利用プランは、発表の「本日提供開始」を超える詳細が確認されていない。エージェントの出力品質やレビュー負担に関する説明も提供元の主張であり、今回の資料にはそれを測定した結果はない。\autocite{src-648195cc00b1,src-a2492a8b244f}
\end{claimboundary}
\begin{claimboundary}
Messages連携の対応機種やWork／Codexでの利用範囲について、二次報道にある詳細は取得した一次資料で別途確認されていない。Appleの推奨・承認を示すものとも扱わない。端末上の自動操作という仕組みの説明も二次報道によるもので、取得した一次資料では確認されていない。\autocite{src-3930e8dede76}
\end{claimboundary}
\begin{claimboundary}
JetBrains向けCopilotの資料の9月8日という日付は取得日であり、ページの発表日ではない。対象週との対応は変更記録の文脈と当時の観測に基づく。変更記録を超える展開時期や利用プランの詳細は確認されていない。\autocite{src-9f289c801bd1}
\end{claimboundary}
\begin{claimboundary}
ChatGPTの操作履歴の基本機能を記した8月14日の記録は日付のみで、対象期間の開始時刻との前後関係を確定できない。この週の変更として扱うのは8月20日の拡張である。プライバシー上の性質は記録の説明を超えて独立検証されていない。\autocite{src-3930e8dede76}
\end{claimboundary}
\begin{claimboundary}
Grok Botの現在の案内ページは編集後の8月26日付であり、8月21日時点の本文を示す資料としては扱わない。8月21日の出来事としての時期は、当時の公式X投稿の観測と同日付の二次資料に基づく。現在のページにある無料試用や企業向け順番待ちの詳細には、8月21日後の編集内容が含まれる可能性がある。\autocite{src-a68f4a2bcd37}
\end{claimboundary}
\begin{claimboundary}
ReplitのFree Modeの対象者や利用上限、GPT-5.6 Lunaの詳しい性能は今回の資料では確認されていない。値下げが無償提供を可能にしたという因果の説明は提供元のものであり、監査によって確かめられた結果ではない。\autocite{src-2f00013902a2}
\end{claimboundary}
\begin{claimboundary}
Google Cloudの記事は8月21日付だが公開時刻がなく、22:00 UTCの締切以前だったかは確定できない。一方、8月20日のAntigravityの記事もバンドル提供を伝えている。Standard／Plus／新興市場向けの条件を超える個別の加入資格については管理画面での確認が必要である。\autocite{src-51b6521f4a58,src-8d408f26c5ec}
\end{claimboundary}
\begin{claimboundary}
Slack向けCopilotの9月8日という日付は資料の取得日であり、対象週との対応は変更記録の文脈と当時の観測に基づく。機能の説明は取得した本文で確認できる範囲に限る。\autocite{src-ab796c22c2d1}
\end{claimboundary}
\begin{claimboundary}
RunwayのMCPによるワークフロー操作について、変更記録には品質や限界を評価する記述がない。機能の追加から品質の向上までは判断しない。\autocite{src-c5fce7aa4ae4}
\end{claimboundary}
\begin{claimboundary}
Teams向けCopilotの9月8日という日付は資料の取得日であり、対象週との対応は変更記録の文脈と当時の観測に基づく。機能の説明は取得した本文で確認できる範囲に限る。\autocite{src-64ccd02222eb}
\end{claimboundary}
\begin{claimboundary}
Kimi Code CLIの新しい道具については機能の説明があるが、振る舞いや品質への効果を示す記述はない。機能の追加からその効果までは判断しない。\autocite{src-9191d43918fc}
\end{claimboundary}
